
\def\PUDcodigo{EME118}

\if\COMPILAR0
   \def\PUDcodacad{04507.62}
\else
   \def\PUDcodacad{04506.34}
\fi

\def\PUDdisciplina{Mecânica dos Fluidos}

\if\COMPILAR0
    \def\PUDsemestre{Opt}
\else
    \def\PUDsemestre{S7}
\fi

\def\PUDhoras{80}

%NOVOS ---------------------------------------------------
\def\PUDhorasTeoria{80}
\def\PUDhorasPratica{0}

\if\COMPILAR0
    \def\PUDinterECA{ \hyperref[PUD:ACE011]{Física 2 (04507.13)} }
\else
    \def\PUDinterECA{ \hyperref[PUD:EME115]{Mecanismos (04506.25)}
    
    \hyperref[PUD:EME114]{Termodinâmica (04506.27)}
    
    \hyperref[PUD:EME119]{Sistemas Mecânicos (04506.35)}
    }
\fi

\def\PUDatuaisECA{---}

\def\PUDrecursos{Projetor multimídia; Quadro branco e pincel.}

%----------------------------------------------------------

\def\PUDcreditos{4}

\if\COMPILAR0
    \def\PUDprerequisito{ \hyperref[PUD:ACE011]{ACE011} }
\else
    \def\PUDprerequisito{ \hyperref[PUD:EME114]{EME114} }
\fi

\if\COMPILAR0
    \def\PUDpreacad{ \hyperref[PUD:ACE011]{Física 2 (04507.13)} }
\else
    \def\PUDpreacad{ \hyperref[PUD:EME114]{Termodinâmica (04506.27)} }
\fi

\def\PUDobjetivos{Conhecer a mecânica dos fluidos clássica, preparando-se para subsequentes estudos em áreas como bombeamento, refrigeração e máquinas térmicas; Preparar-se para o uso efetivo da teoria da mecânica dos fluidos na prática da engenharia;  Compreender os fenômenos relativos à conservação da massa, à conservação da quantidade de movimento, à conservação da energia.}

\def\PUDementa{Propriedades dos fluidos; Estática dos fluidos; Cinemática dos fluidos; Equação da energia para regime permanente; Equação da quantidade de movimento para regime permanente; Análise dimensional - semelhança; Escoamento permanente de fluido incompressível em condutos forçados; Noções de instrumentação para medida das propriedades dos fluidos; Fluidodinâmica.}

\def\PUDconteudo{ & Introdução, definição e propriedades dos fluidos: Conceitos fundamentais e definição de fluido; Tensão de cisalhamento – Lei de Newton da Viscosidade;  Viscosidade absoluta ou dinâmica;  Simplificação prática;  Massa específica;  Peso específico;  Peso específico relativo para líquidos;  Viscosidade cinemática;  Fluido ideal;  Fluido ou escoamento incompressível;  Equação de estado dos gases;  (8 horas-aulas). 
\\ 
 & Estática dos fluidos: Pressão;  Teorema de Stevin;  Pressão em torno de um ponto de um fluido em repouso;  Lei de Pascal;  Carga de pressão;  Escalas de pressão;  Unidades de pressão;  O barômetro;  Medidores de pressão;  Força numa superfície plana submersa;  Centro das pressões;  Força em superfícies reversas, submersas;  Empuxo;  Flutuador - Nomenclatura;  Estabilidade;  Estabilidade vertical;  Estabilidade à rotação;  Equilíbrio relativo;  Recipiente com movimento de translação uniformemente acelerado segundo a horizontal;  Recipiente com movimento de translação uniformemente acelerado segundo a vertical;  Recipiente com movimento de translação uniformemente acelerado ao longo de um plano inclinado;  Recipiente com movimento de rotação de velocidade angular constante;  (10 horas-aulas).  
\\ 
 & Cinemática dos fluidos: Regimes ou movimentos variado e permanente;  Escoamento laminar e turbulento;  Trajetória e linha de corrente;  Escoamento unidimensional ou uniforme na seção;  Vazão – velocidade média na seção;  Equação da continuidade para o regime permanente;  Velocidade e aceleração nos escoamentos de fluidos (6 horas-aulas). 
\\ 
 & Equação da energia para regime permanente: Tipos de energias mecânicas associadas a um fluido;  Equação de Bernoulli;  Equação da energia e presença de uma máquina;  Potência da máquina e noção de rendimento;  Equação da energia para fluido real;  Diagrama de velocidades não-uniforme na seção;  Equação da energia para diversas entradas e saídas e escoamento em regime permanente de um fluido incompressível, sem trocas de calor;  Interpretação da perda de carga;  Equação da energia geral para regime permanente.  (8 horas-aulas).  
\\ 
 & Equação da quantidade de movimento para regime permanente: Equação da quantidade de movimento;  Método de utilização da equação;  Forças em superfícies sólidas em movimento;  Equação da quantidade de movimento para diversas entradas e saídas em regime permanente.  (6 horas-aulas). 
\\ 
 & Análise dimensional - semelhança: Grandezas fundamentais e derivadas.  Equações dimensionais;  Sistemas coerentes de unidades;  Números adimensionais;  Vantagem da utilização dos números adimensionais na pesquisa de uma lei física;  Teorema dos PI;  Alguns números adimensionais típicos;  Semelhança ou teoria dos modelos;  Escalas de semelhança;  relações entre escalas.  (6 horas-aulas).  
\\ 
 & Escoamento permanente de fluido incompressível em condutos forçados:.  Definições;  Estudo da perda de carga distribuída;  Fórmula da perda de carga distribuída;  Experiência de Nikuradse;  Condutos industriais;  Problemas típicos envolvendo apenas perda de carga distribuída;  Perdas de carga singulares;  Instalações de recalque;  Linhas de energia e piezométrica (10 horas-aulas). 
\\ 
 & Noções de instrumentação para medida das propriedades dos fluidos: Massa específica e peso específico relativo;  Viscosidade;  Medida da velocidade com tubo de Pitot;  Medida da vazão.  (6 horas-aulas). 
\\ 
 & Fluidodinâmica: Conceitos fundamentais;  Força de arrasto de superfície;  Força de arrasto de forma ou de pressão;  Força de arrasto total;  Força de sustentação;  Máquinas de Fluxo (8 horas-aulas). }

\def\PUDmetodo{Aulas expositórias teóricas, sendo também utilizada a metodologia de ensino baseada em projetos de sistemas de bombeamentos.}

\def\PUDavalia{A avaliação será feita com aplicação de provas teóricas, além da possibilidade de inclusão de trabalhos e seminários no decorrer da disciplina. Uma das notas da disciplina é dada ao projeto do sistema de bombeamento com memorial de cálculo.}

\def\PUDbibbasica{ & \href{www.biblioteca.ifce.edu.br/index.asp?codigo_sophia=63133}{Fox}, R. W.; MacDolnad, A. T.; Pritchard, P. J.  Introdução à Mecânica dos Fluidos.  8ª ed.  Rio de Janeiro, RJ: LTC, 2014.  871p.  ISBN: 9788521623021. 
\\ 
 &  \href{www.biblioteca.ifce.edu.br/index.asp?codigo_sophia=24162}{Brunetti}, F.  Mecânica dos Fluidos.  2ª ed.  São Paulo, SP: Pearson Prentice Hall, 2008.  431p.  ISBN: 9788576051824. 
\\ 
 & \href{www.biblioteca.ifce.edu.br/index.asp?codigo_sophia=23814}{White}, Frank M.  Mecânica dos fluidos.  6º ed.  Porto Alegre, RS: AMGH, 2011.  880p.  ISBN: 9788563308214.
 }

\def\PUDbibcomplementar{ & \href{www.biblioteca.ifce.edu.br/index.asp?codigo_sophia=63116}{Post}, S.  Mecânica dos Fluidos Aplicada e Computacional.  1ª ed.  Rio de Janeiro, RJ: LTC, 2013.  402p.  ISBN: 9788521620990.
\\ 
 & \href{www.biblioteca.ifce.edu.br/index.asp?codigo_sophia=65410}{Borgnakke}, C.; Sonntag, R. E.  Fundamentos da Termodinâmica.  São Paulo, SP: Edgard Blücher, 2013.  728p.  ISBN: 9788521207924. 
\\ 
 &  \href{www.biblioteca.ifce.edu.br/index.asp?codigo_sophia=24058}{Moran}, M. J.; Shapiro, H. N.; Muson, B. R.; DeWitt, D. P.  Introdução à engenharia de sistemas térmicos: termodinâmica, mecânica dos fluidos e transferência de calor.  Rio de Janeiro, RJ : LTC, 2005.  604p.  ISBN: 8521614463.
\\
 &  \href{www.biblioteca.ifce.edu.br/index.asp?codigo_sophia=41293}{Baptista}, Márcio.  Fundamentos de engenharia hidráulica.  3º ed.  Belo Horizonte, MG: UFMG, 2014.  473p.  ISBN: 9788570418289.
\\
 &  \href{www.biblioteca.ifce.edu.br/index.asp?codigo_sophia=41951}{Macintyre}, A. J.  Instalações hidráulicas: prediais e industriais.  4ª ed.  Rio de Janeiro, RJ: LTC, 2010.  579p.  ISBN: 9788521616573.
 }

\def\PUDautor{Francisco Frederico dos Santos Matos}

\def\PUDdata{2014-04-23}

\def\PUDrevisao{3}

\def\PUDrevisor{Francisco Frederico dos Santos Matos}

\def\PUDdatarevisao{2019-05-15}

\PUDcorpo
